% !PRINTED 536
\weberpage{536}
%<<CONT
fixed. If then $\pi_i$ is a permutation that carries the figure 0 into $i$, then $\pi_i^{-1}\,Q\,\pi_i$ is the group of the permutations in $P$ that do not alter the figure $i$. But since in every composition of permutations from $P$ the components may be interchanged, we have
\[
\pi_i^{-1}\,Q\,\pi_i = \pi_i^{-1}\pi_i\,Q = Q,
\]
that is, the group $Q$ also leaves the figure $i$ unchanged. On account of the transitivity of $P$, however, $i$ may denote any of the figures $1,2\ldots m-1$; consequently $Q$ consists of the single identical permutation, and in $P$ there is, besides the identical one, no permutation that leaves a figure unchanged. But if one adjoins \gesperrt{a} root $\alpha$, the group $P$ reduces to the unit group, or the equation is solved; $\varphi(x)=0$ is a normal equation and hence its own \gesperrt{Galois} resolvent. If $\alpha_k=\Theta_k(\alpha)$, then the \gesperrt{Galois} group of this equation consists of the substitutions:
\[
\sigma_k = [\alpha,\Theta_k(\alpha)],
\]
and we have
\[
% [sic] authority prints x in the second composite; the parallel expression and the next line require \alpha.
\sigma_h\sigma_k = [\alpha,\Theta_h\Theta_k(\alpha)],\quad \sigma_k\sigma_h = [\alpha,\Theta_k\Theta_h(x)].
\]
Since the group is to be commutative, we must have
\[
\Theta_h\Theta_k(\alpha) = \Theta_k\Theta_h(\alpha)
\]
that is, $\varphi(x)=0$ is an \gesperrt{Abelian equation}. This is the reason why the commutative groups are also called \gesperrt{Abelian} groups (\S.\,148). It therefore also follows from the theorem given above on the group of normal equations that in a transitive \gesperrt{Abelian} permutation group the number of permutations agrees with the number of the permuted figures.

Here we have presupposed the irreducibility of the given equation. If we wish to take into consideration also reducible equations with commutative group, we need only put in place of the equation itself one of its \gesperrt{Galois} resolvents.

It remains, finally, to remark that in a commutative group every divisor is normal, since we always have $\pi^{-1}\varkappa\pi = \varkappa$ when $\pi$ and $\varkappa$ are arbitrary elements of a commutative group.
